% BPCI / BPI / Herm�s Proof Log
% This LaTeX document is a structured proof log for real demos and tests.

% NOTE: The original standalone \documentclass has been commented out so
% that this file can be included from main_pravyom_report.tex.
% If you want to compile this file standalone again, uncomment and
% add \begin{document}/\end{document} wrappers.
% \documentclass[11pt,a4paper]{report}
\usepackage[margin=1in]{geometry}
\usepackage[T1]{fontenc}
\usepackage[utf8]{inputenc}
\usepackage{hyperref}
\usepackage{courier}
\usepackage{listings}
\usepackage{longtable}

\lstset{
  basicstyle=\footnotesize\ttfamily,
  breaklines=true,
  columns=fullflexible,
}

% \begin{document}
% This file is intended to be used as Chapter 2 of a larger thesis/report.
% If included in a main document, remove the \documentclass line above and
% the \begin{document}/\end{document} wrapper, and input this file instead.

% Adjust the chapter counter so this becomes Chapter 2 in the final document.
% Comment out the next line if you want it to be Chapter 1 instead.
\setcounter{chapter}{1}

\chapter{BPCI / BPI / Herm\`es Proof Log}

\section{Purpose}

This document collects \textbf{real execution logs} and \textbf{demo reports} from the BPCI constellation, BPI\,$\leftrightarrow$\,BPCI lifecycle, and Herms P2P mesh. It is intended as a \emph{proof log} that can be attached to architecture and readiness reports.

Each section includes:

\begin{itemize}
  \item The exact command(s) used to run the demo or test.
  \item The path where the report/log file is written (usually under \/tmp).
  \item A short narrative of what the demo proves.
  \item A placeholder where the full log/report can be pasted or attached.
\end{itemize}

\section{Environment Summary}

Fill in the concrete values for the run you are documenting:

\begin{itemize}
  \item Host OS: Linux (e.g. Ubuntu 22.04), kernel version: \texttt{<uname -r>}.
  \item Rust toolchain: \texttt{rustc --version}, \texttt{cargo --version}.
  \item Repository root: \texttt{/home/umesh/metanode}.
  \item Crates: \texttt{bpi-core}, \texttt{bpci-enterprise} (a.k.a. \texttt{pravyom-enterprise}).
  \item Build command: \texttt{cargo build -p bpi-core -p bpci-enterprise} (debug profile).
\end{itemize}

Optionally include the exact command transcript from the build and selected tests in an appendix.

\section{BPCI Constellation 1-Minute Demo}

\subsection{Constellation Setup}

The BPCI constellation consists of:

\begin{itemize}
  \item Component 1: LCCD Consensus Server (HTTP, port 9001).
  \item Component 2: BPCI Blockchain Server (HTTP API, port 8082).
  \item Component 3: BPCI Auction Mempool Server (HTTP, port 9004).
  \item Component 4: BPCI Auction DB Server (local maintainer, HTTP, port 7002).
  \item Component 5: BPCI--BPI Bridge (HTTP, port 6001).
\end{itemize}

\paragraph{Controller binary (layout and commands).}
From repo root:

\begin{lstlisting}[language=bash]
cargo run -p pravyom-enterprise --bin bpci_constellation_control
\end{lstlisting}

\paragraph{Environment bindings used by the demo.}

\begin{lstlisting}[language=bash]
export DEPLOYMENT_TYPE="BSO-K8 orchestrator"
export NETWORK_BINDING="0.0.0.0 (external access)"
export CLUSTER_NAME="bpci-local-dev"
export NAMESPACE="bpci-enterprise"

export BPCI_CONSENSUS_URL="http://127.0.0.1:9001"
export BPCI_BLOCKCHAIN_URL="http://127.0.0.1:8082"
export BPCI_AUCTION_MEMPOOL_URL="http://127.0.0.1:9004"
export BPCI_AUCTION_DB_URL="http://127.0.0.1:7002"
\end{lstlisting}

\paragraph{Component startup commands.}
Each command is typically run in its own terminal.

\begin{lstlisting}[language=bash]
# Component 1 - LCCD Consensus Server
cargo run -p pravyom-enterprise --bin bpci-consensus-server \
  -- --port 9001

# Component 2 - BPCI Blockchain Server
cargo run -p pravyom-enterprise --bin bpci_blockchain_server \
  -- --api-port 8082 \
     --consensus-server-url http://127.0.0.1:9001

# Component 3 - BPCI Auction Mempool Server
cargo run -p pravyom-enterprise --bin bpci_auction_mempool_server \
  -- --api-port 9004

# Component 4 - BPCI Auction DB Server (local)
cargo run -p pravyom-enterprise --bin bpci_auction_db_server

# Component 5 - BPCI-BPI Bridge (HTTP on 6001)
cargo run -p pravyom-enterprise --bin bpci_bpi_bridge \
  -- --port 6001
\end{lstlisting}

\paragraph{Health checks before running the demo.}

\begin{lstlisting}[language=bash]
curl -i "$BPCI_CONSENSUS_URL/api/v1/health"
curl -i "$BPCI_BLOCKCHAIN_URL/health"

curl -i "$BPCI_AUCTION_MEMPOOL_URL/auction/submit" \
  -X POST -H 'Content-Type: application/json' -d '{}'

curl -i "$BPCI_AUCTION_DB_URL/api/v1/auction/record" \
  -X POST -H 'Content-Type: application/json' \
  -d '{"tx_id":"test","from_bpi":"bpi","to_bpci":"bpci","amount":1,"record_type":"demo"}'
\end{lstlisting}

\subsection{Demo Command and Report}

\paragraph{Demo run.}

\begin{lstlisting}[language=bash]
# From /home/umesh/metanode
cargo run -p pravyom-enterprise --bin bpci_constellation_demo
\end{lstlisting}

The demo writes a human-readable report to:

\begin{itemize}
  \item \texttt{/tmp/bpci_constellation_demo_report.txt}
\end{itemize}

\paragraph{What this demo proves.}

\begin{itemize}
  \item Real LCCD consensus health endpoint responds (HTTP 200) during the run.
  \item LCCD consensus rounds can be started per transaction (\texttt{/api/v1/lccd/consensus/start}).
  \item Blockchain server accepts and records synthetic transactions via \texttt{/api/v1/transactions}.
  \item Auction mempool server accepts submissions via \texttt{/auction/submit}.
  \item Auction DB maintainer records entries via \texttt{/api/v1/auction/record}.
  \item The per-transaction table shows integrated behaviour across all 4--5 components.
\end{itemize}

\paragraph{Captured report (paste full contents).}

Paste the full contents of \texttt{/tmp/bpci_constellation_demo_report.txt} below, e.g. using a verbatim or listings environment:

\begin{lstlisting}
BPCI Constellation 1-Minute Demo Report
=======================================

Account: bpi-constellation-demo-001
Initial balance: 200 BPI
Final balance  : 139 BPI
Duration (s)   : 16
Total tx       : 8

idx | amount | fee | total | status | cons_hlth | lccd_started | lccd_status | lccd_round | blockchain | auction | db
----+--------+-----+-------+--------+-----------+-------------+-------------+------------+-----------+--------+--------
  0 |      5 |   1 |     6 | accepted | healthy     | yes         | started     | cc55546f-… | submitted | submit… | recorded
  1 |      6 |   1 |     7 | accepted | healthy     | yes         | started     | 84c1f605-… | submitted | submit… | recorded
  2 |      7 |   1 |     8 | accepted | healthy     | yes         | started     | 8fc8c34a-… | submitted | submit… | recorded
  3 |      8 |   1 |     9 | accepted | healthy     | yes         | started     | 05944216-… | submitted | submit… | recorded
  4 |      9 |   1 |    10 | accepted | healthy     | yes         | started     | 26fec299-… | submitted | submit… | recorded
  5 |      5 |   1 |     6 | accepted | healthy     | yes         | started     | 9fac74ac-… | submitted | submit… | recorded
  6 |      6 |   1 |     7 | accepted | healthy     | yes         | started     | 93037102-… | submitted | submit… | recorded
  7 |      7 |   1 |     8 | accepted | healthy     | yes         | started     | e8dad90a-… | submitted | submit… | recorded
\end{lstlisting}

\section{BPI \texorpdfstring{$\leftrightarrow$}{<->} BPCI Lifecycle Demo}

\subsection{Command and Report}

\begin{lstlisting}[language=bash]
# From /home/umesh/metanode
cargo run -p pravyom-enterprise --bin bpi_bpci_lifecycle_demo
\end{lstlisting}

Report path:

\begin{itemize}
  \item \texttt{/tmp/bpi_bpci_lifecycle_demo_report.txt}
\end{itemize}

\paragraph{What this demo proves.}

\begin{itemize}
  \item BPI-side accounting of principal + gas fees across a sequence of demo transactions.
  \item How BPI balances move when BPCI infrastructure is treated as an external service.
  \item Time-series economic data that can be compared to BPCI-side constellation behaviour.
\end{itemize}

\paragraph{Captured report (paste full contents).}

\begin{lstlisting}[language=text]
BPI→BPCI Testnet Lifecycle Mini-Demo Report
==========================================

1. Overview
-----------
Account address    : bpi-testnet-user-001
Plan name          : Testnet
Plan monthly cost  : 10.00 CAD
Monthly allocation : 1000 BPI
Free allocation    : 200 BPI
Gas fee percentage : 0.500%
Start time (UTC)   : 2025-11-25T20:03:05.532955122+00:00
End time (UTC)     : 2025-11-25T20:03:05.533060608+00:00
Duration (seconds) : 0
Total tx attempted : 20
Tx accepted        : 20
Tx rejected        : 0

2. Account Lifecycle Summary
----------------------------
Initial free balance : 200 BPI
Final balance        : 40 BPI
Total principal sent : 140 BPI
Total gas fees       : 20 BPI
Total usage          : 160 BPI

3. BPI Bundle Flow Timeline (offline-simulated)
-----------------------------------------------
idx | timestamp                  | amount | fee | total | bal_before | bal_after | status                         | bundle_id
----+----------------------------+--------+-----+-------+-----------+-----------+------------------------------+----------
  0 | 2025-11-25T20:03:05Z |      5 |   1 |     6 |       200 |       194 | accepted                     | bundle_1
  1 | 2025-11-25T20:03:05Z |      6 |   1 |     7 |       194 |       187 | accepted                     | bundle_1
  2 | 2025-11-25T20:03:05Z |      7 |   1 |     8 |       187 |       179 | accepted                     | bundle_1
  3 | 2025-11-25T20:03:05Z |      8 |   1 |     9 |       179 |       170 | accepted                     | bundle_1
  4 | 2025-11-25T20:03:05Z |      9 |   1 |    10 |       170 |       160 | accepted                     | bundle_1
  5 | 2025-11-25T20:03:05Z |      5 |   1 |     6 |       160 |       154 | accepted                     | bundle_2
  6 | 2025-11-25T20:03:05Z |      6 |   1 |     7 |       154 |       147 | accepted                     | bundle_2
  7 | 2025-11-25T20:03:05Z |      7 |   1 |     8 |       147 |       139 | accepted                     | bundle_2
  8 | 2025-11-25T20:03:05Z |      8 |   1 |     9 |       139 |       130 | accepted                     | bundle_2
  9 | 2025-11-25T20:03:05Z |      9 |   1 |    10 |       130 |       120 | accepted                     | bundle_2
 10 | 2025-11-25T20:03:05Z |      5 |   1 |     6 |       120 |       114 | accepted                     | bundle_3
 11 | 2025-11-25T20:03:05Z |      6 |   1 |     7 |       114 |       107 | accepted                     | bundle_3
 12 | 2025-11-25T20:03:05Z |      7 |   1 |     8 |       107 |        99 | accepted                     | bundle_3
 13 | 2025-11-25T20:03:05Z |      8 |   1 |     9 |        99 |        90 | accepted                     | bundle_3
 14 | 2025-11-25T20:03:05Z |      9 |   1 |    10 |        90 |        80 | accepted                     | bundle_3
 15 | 2025-11-25T20:03:05Z |      5 |   1 |     6 |        80 |        74 | accepted                     | bundle_4
 16 | 2025-11-25T20:03:05Z |      6 |   1 |     7 |        74 |        67 | accepted                     | bundle_4
 17 | 2025-11-25T20:03:05Z |      7 |   1 |     8 |        67 |        59 | accepted                     | bundle_4
 18 | 2025-11-25T20:03:05Z |      8 |   1 |     9 |        59 |        50 | accepted                     | bundle_4
 19 | 2025-11-25T20:03:05Z |      9 |   1 |    10 |        50 |        40 | accepted                     | bundle_4

4. Bundle & Auction Summary (Conceptual)
----------------------------------------
Total bundles formed : 4
- bundle_1 → tx_count=5 principal=35 gas_fees=5
- bundle_2 → tx_count=5 principal=35 gas_fees=5
- bundle_3 → tx_count=5 principal=35 gas_fees=5
- bundle_4 → tx_count=5 principal=35 gas_fees=5

5. Component Lifecycle (Conceptual)
-----------------------------------
- Component 1 (Consensus): In this offline demo, consensus readiness is assumed. In production,
  the bridge would call the consensus health endpoint before processing a bundle.
- Component 2 (Blockchain): Transaction submissions and fee accounting would be recorded via HTTP.
  Here we simulate only the accounting side (principal + gas fee debits).
- Component 3 (Auction Mempool): In production, each tx becomes a candidate in the auction mempool.
  This demo models the BPI-side bundle economics without a real mempool.
- Component 4 (Auction DB): Persistent records would be written via the DB maintainer. Here we
  capture a time-series table instead, suitable for testnet analysis.

5. Observations
----------------
- This demo is fully in-memory and offline; it does not contact any external servers.
- Gas fee calculation matches the bridge testnet plan: 0.500% with a minimum 1 BPI per tx.
- The report shows the full BPI token lifecycle for a single testnet account: from free allocation,
  through multiple bundle-like sends, to final remaining balance and total usage.
- It is designed to be laptop-safe and quick to run, ideal for reviewers and testnet tuning.
\end{lstlisting}

\section{Hermes-Lite Web-4 Mesh Demo}

\subsection{Command and Report}

\begin{lstlisting}[language=bash]
# From /home/umesh/metanode
cargo run -p pravyom-enterprise --bin hermes_mesh_demo
\end{lstlisting}

Report path:

\begin{itemize}
  \item \texttt{/tmp/hermes_mesh_demo_report.txt}
\end{itemize}

\paragraph{What this demo proves.}

\begin{itemize}
  \item Hermes-Lite mesh can be instantiated in-memory on a laptop.
  \item Consensus rounds propagate through a small 3-node mesh (local + 2 bootstrap nodes).
  \item Mesh health ratio, kappa metrics, and throughput evolve over time.
  \item A single cellular division event is exercised mid-run.
\end{itemize}

\paragraph{Captured report (paste full contents).}

\begin{lstlisting}[language=text]
% --- BEGIN hermes_mesh_demo_report.txt ---
% (paste the entire report here)
% --- END hermes_mesh_demo_report.txt ---
\end{lstlisting}

\section{Hermes P2P Decentralization Phase Simulation}

\subsection{Command and Report}

\begin{lstlisting}[language=bash]
# From /home/umesh/metanode
cargo run -p pravyom-enterprise --bin hermes_decentralization_demo
\end{lstlisting}

Report path:

\begin{itemize}
  \item \texttt{/tmp/hermes_decentralization_demo_report.txt}
\end{itemize}

\paragraph{What this demo proves.}

\begin{itemize}
  \item How BPI OS (BPIOS) nodes are gradually attached to the Hermes mesh over time.
  \item How multiple cellular division events increase the number of consensus "cells" in the mesh.
  \item How the system transitions through phases:
    \begin{itemize}
      \item Centralized Constellation
      \item NodeConnectionSync (Hermes P2P forming)
      \item MeshEvolution (decentralization in progress)
      \item AutonomousMesh (full decentralization)
    \end{itemize}
  \item At which round / node count the AutonomousMesh threshold (nodes, health ratio, divisions) is reached, if at all.
\end{itemize}

\paragraph{Captured report (paste full contents).}

\begin{lstlisting}[language=text]
% --- BEGIN hermes_decentralization_demo_report.txt ---
% (paste the entire report here)
% --- END hermes_decentralization_demo_report.txt ---
\end{lstlisting}

\section{Additional Tests and Logs}

Use this section to add any further tests from the BPCI-side test plan (e.g. cluster ledger, readiness, security, or economics tests). For each test or suite, follow the same pattern:

\begin{enumerate}
  \item Document the exact command(s) used.
  \item State clearly what behaviour or property the test demonstrates.
  \item Paste the key parts of the log or point to an attached file.
\end{enumerate}

\subsection{Example Placeholder}

\paragraph{Command.}

\begin{lstlisting}[language=bash]
# Example placeholder for a future test
# cargo run -p pravyom-enterprise --bin <binary> -- <flags>
\end{lstlisting}

\paragraph{What this test proves.}

\begin{itemize}
  \item Describe the invariant or readiness property being exercised.
\end{itemize}

\paragraph{Captured log or report.}

\begin{lstlisting}[language=text]
% --- BEGIN example_log.txt ---
% (paste log excerpt here)
% --- END example_log.txt ---
\end{lstlisting}

% \end{document}
